\documentclass{article}
\usepackage{mathtools}
\usepackage{graphicx}
\usepackage{xcolor}
\usepackage{url}
\usepackage[margin=1in]{geometry}
\usepackage{filecontents}


\begin{document}
\title{Peer Review - Aaron \\
Augmenting Decompiler Output with Learned Variable Names and Types}
\author{Dustin Randall}
\date{}
\maketitle

\section{What was interesting}
This topic was interesting to me because I have looked at compiled C++ code and realize that it is so very different from the original code.
Decompiling binaries back into readable source code is used for a variety of applications, including by IDEs when exceptions are thrown from compiled libraries.
Seeing approaches for augmenting decompiler output seems like a very challenging space, and I wonder if the approach may be overfit to existing C code.

\section{Why is it important}
The ability to decompile binaries is also used in reverse engineering and automated program repair.
Improving the quality of the decompiled output can improve the ability to understand and eventually modify the original binary using binary patching.
Investigating techniques for augmenting decompiler output can also contribute to diagnosing code gen errors from compilers which can accelerate new language feature development.

\section{Presentation}
I thought the visuals used were well done and helped in understanding the topics.
I generally am not a fan of people sitting while presenting, but it wasn't overly distracting.
One potential additional visual would be showing a side-by-side comparison of surprising compiler optimizations, potentially using godbolt.org.

\section{Citations}
Beyond {Classification}: {Inferring} {Function} {Names} in {Stripped} {Binaries} via {Domain} {Adapted} {LLMs}~\cite{jiang_beyond_2025}
and Decompile-{Bench}: {Million}-{Scale} {Binary}-{Source} {Function} {Pairs} for {Real}-{World} {Binary} {Decompilation}~\cite{tan_decompile-bench_2025}

\bibliography{citations}
\bibliographystyle{plain}

\end{document}